\begin{proof}[Proof of \cref{obj:lem:variance-cost-envelope}]
\begin{enumerate}
\item We use the classical polynomial identities supplied by
\cref{obj:lem:classical-legendre-facts}. The factorial-series estimate in
\cref{obj:lem:exact-covariance} supplies an absolute constant
\[
C_{\mathrm{fac}}\in(0,\infty)
\]
such that, for every integer \(L\ge1\) and every \(\sigma\in[0,1]\),
\[
m_L=3(L-1),\qquad
V_L(\sigma)\le C_{\mathrm{fac}}L^2
\sum_{j=0}^{m_L}
\frac{(C_{\mathrm{fac}}\sigma^2L^4)^j}{(j!)^3}.
\]
Throughout the proof, write
\[
h_L=L^{-2}>0.
\]
% lean: CausalSmith.Stat.RdTruesideNoiseFrontier.exact_covariance

\item First construct an envelope for the first two noise regimes. Define
\[
C_{\mathrm{nc}}
=\max\{C_{\mathrm{fac}},\,3(C_{\mathrm{fac}}+1)\}>0,
\qquad
\eta_L(\sigma)=\sigma L^2\ge0,
\qquad
\zeta_L(\sigma)=(C_{\mathrm{fac}}+1)\eta_L(\sigma)^{2/3}\ge0.
\]
These definitions apply to every \(L\ge1\) and \(\sigma\in[0,1]\). Since
\[
\bigl(\eta_L(\sigma)^{2/3}\bigr)^3=\eta_L(\sigma)^2,
\qquad
C_{\mathrm{fac}}\le(C_{\mathrm{fac}}+1)^3,
\]
we have
\[
C_{\mathrm{fac}}\sigma^2L^4
\le \zeta_L(\sigma)^3.
\]

For nonnegative summands, the sum of their cubes is bounded by the cube
of their sum. This follows by induction from
\[
(\xi+\omega)^3
=\xi^3+\omega^3+3\xi\omega(\xi+\omega)
\ge \xi^3+\omega^3,
\qquad \xi,\omega\in[0,\infty).
\]
Applying this inequality and then the nonnegative exponential series gives
\[
\begin{aligned}
\sum_{j=0}^{m_L}
\frac{(C_{\mathrm{fac}}\sigma^2L^4)^j}{(j!)^3}
&\le
\sum_{j=0}^{m_L}
\left(\frac{\zeta_L(\sigma)^j}{j!}\right)^3\\
&\le
\left(\sum_{j=0}^{m_L}\frac{\zeta_L(\sigma)^j}{j!}\right)^3\\
&\le \exp\{3\zeta_L(\sigma)\}\\
&\le \exp\{C_{\mathrm{nc}}\eta_L(\sigma)^{2/3}\}.
\end{aligned}
\]
Consequently,
\[
V_L(\sigma)
\le C_{\mathrm{nc}}L^2
\exp\{C_{\mathrm{nc}}\eta_L(\sigma)^{2/3}\}.
\]

If \(\sigma\le h_L\), then
\[
0\le\eta_L(\sigma)\le1,
\qquad
\eta_L(\sigma)^{2/3}\le1
=1+\Psi(h_L,\sigma).
\]
If \(\sigma>h_L\) and \(\sigma^4\le h_L\), the intermediate branch instead gives
\[
1+\Psi(h_L,\sigma)
=(\sigma/h_L)^{2/3}
=\eta_L(\sigma)^{2/3}.
\]
Thus the same constant satisfies
\[
V_L(\sigma)\le C_{\mathrm{nc}}L^2
\exp\{C_{\mathrm{nc}}[1+\Psi(h_L,\sigma)]\}
\quad\text{whenever}\quad
\sigma\le h_L\ \text{or}\ \sigma^4\le h_L.
\]
% lean: CausalSmith.Stat.RdTruesideNoiseFrontier.kernelVariance_noncompact_noise_envelope

\item Next construct the compact-branch envelope. Define the absolute constant
\[
C_{\mathrm{comp}}
=\max\{C_{\mathrm{fac}},\,7+3\log(C_{\mathrm{fac}}+1)\}>0.
\]
Suppose that \(\sigma>h_L\) and \(\sigma^4>h_L\), and define
\[
\eta_{\mathrm{comp}}=\sigma^2L,
\qquad
\alpha_{\mathrm{comp}}=(C_{\mathrm{fac}}+1)\eta_{\mathrm{comp}}.
\]
Multiplication of \(\sigma^4>L^{-2}\) by \(L^2>0\) yields
\[
\eta_{\mathrm{comp}}^2=\sigma^4L^2>1.
\]
Since \(\eta_{\mathrm{comp}}\ge0\), it follows that
\[
\eta_{\mathrm{comp}}\ge1,
\qquad
\alpha_{\mathrm{comp}}\ge1.
\]

We first establish the finite-series estimate used here. For
\[
\alpha\in[1,\infty),
\qquad
N_{\mathrm{cut}}\in\mathbb Z_{\ge0},
\qquad
N_{\mathrm{cut}}\le3L+1,
\]
every integer \(0\le j<N_{\mathrm{cut}}\) satisfies \(j\le3L\).
The nonnegative exponential series gives \(L^j/j!\le e^L\), and therefore
\[
\frac{(\alpha L^3)^j}{(j!)^3}
=\alpha^j\left(\frac{L^j}{j!}\right)^3
\le \alpha^{3L}e^{3L}.
\]
The number of summands satisfies
\[
N_{\mathrm{cut}}\le3L+1\le4L+1\le e^{4L},
\]
where the last inequality is \(1+u\le e^u\) at \(u=4L\). Hence
\[
\sum_{j=0}^{N_{\mathrm{cut}}-1}
\frac{(\alpha L^3)^j}{(j!)^3}
\le N_{\mathrm{cut}}\alpha^{3L}e^{3L}
\le \exp\{L(7+3\log\alpha)\}.
\]
For \(N_{\mathrm{cut}}=0\), the sum is empty and equals zero, so this estimate
also covers that case.

Now \(m_L+1\le3L+1\), and
\[
C_{\mathrm{fac}}\sigma^2L^4
\le \alpha_{\mathrm{comp}}L^3.
\]
Applying the preceding estimate gives
\[
\sum_{j=0}^{m_L}
\frac{(C_{\mathrm{fac}}\sigma^2L^4)^j}{(j!)^3}
\le
\exp\{L(7+3\log\alpha_{\mathrm{comp}})\}.
\]
On this branch,
\[
\sqrt{h_L}=L^{-1},
\qquad
h_L^{-1/2}=L,
\qquad
1+\Psi(h_L,\sigma)
=L\bigl(1+\log\eta_{\mathrm{comp}}\bigr).
\]
Moreover,
\[
\log\alpha_{\mathrm{comp}}
=\log(C_{\mathrm{fac}}+1)+\log\eta_{\mathrm{comp}},
\qquad
\log\eta_{\mathrm{comp}}\ge0.
\]
Since \(\log(C_{\mathrm{fac}}+1)\ge0\), the definition of
\(C_{\mathrm{comp}}\) implies both
\[
7+3\log(C_{\mathrm{fac}}+1)\le C_{\mathrm{comp}},
\qquad
3\le C_{\mathrm{comp}}.
\]
It follows that
\[
\begin{aligned}
L(7+3\log\alpha_{\mathrm{comp}})
&=L\bigl(7+3\log(C_{\mathrm{fac}}+1)
            +3\log\eta_{\mathrm{comp}}\bigr)\\
&\le C_{\mathrm{comp}}L(1+\log\eta_{\mathrm{comp}})\\
&=C_{\mathrm{comp}}[1+\Psi(h_L,\sigma)].
\end{aligned}
\]
Combining these inequalities with the factorial-series estimate and
\(C_{\mathrm{fac}}\le C_{\mathrm{comp}}\) proves
\[
V_L(\sigma)\le C_{\mathrm{comp}}L^2
\exp\{C_{\mathrm{comp}}[1+\Psi(h_L,\sigma)]\}
\quad\text{if}\quad
\sigma>h_L,\ \sigma^4>h_L.
\]
% lean: CausalSmith.Stat.RdTruesideNoiseFrontier.kernelVariance_compact_noise_envelope

\item Choose the common absolute constant
\[
C=\max\{C_{\mathrm{nc}},C_{\mathrm{comp}}\}>0.
\]
For every \(L\ge1\) and \(\sigma\in[0,1]\), the branch identities above show
\[
1+\Psi(h_L,\sigma)
=
\begin{cases}
1,&\sigma\le h_L,\\
(\sigma L^2)^{2/3},&\sigma>h_L,\ \sigma^4\le h_L,\\
L[1+\log(\sigma^2L)],&\sigma>h_L,\ \sigma^4>h_L
\end{cases}
\ \ge0.
\]
In the last branch, nonnegativity follows from \(\sigma^2L\ge1\).
Thus, for every auxiliary constant
\[
C_{\mathrm{aux}}\in(0,C],
\]
we have
\[
C_{\mathrm{aux}}L^2\le CL^2,
\qquad
C_{\mathrm{aux}}[1+\Psi(h_L,\sigma)]
\le C[1+\Psi(h_L,\sigma)].
\]
Monotonicity and positivity of the exponential therefore yield
\[
C_{\mathrm{aux}}L^2
\exp\{C_{\mathrm{aux}}[1+\Psi(h_L,\sigma)]\}
\le
CL^2\exp\{C[1+\Psi(h_L,\sigma)]\}.
\]
% lean: CausalSmith.Stat.RdTruesideNoiseFrontier.endpoint_noiseCost_factor_nonneg

\item Fix \(L\ge1\) and \(\sigma\in[0,1]\). If \(\sigma\le h_L\), apply the
first envelope with \(C_{\mathrm{aux}}=C_{\mathrm{nc}}\).
If \(\sigma>h_L\), split further according to whether
\(\sigma^4\le h_L\). In that case, apply the same envelope; when
\(\sigma^4>h_L\), apply the compact-branch envelope with
\(C_{\mathrm{aux}}=C_{\mathrm{comp}}\). The enlargement inequality in each
case gives
\[
V_L(\sigma)\le CL^2
\exp\{C[1+\Psi(L^{-2},\sigma)]\}.
\]
The first case includes \(\sigma=0\) and all \(\sigma\in[0,1]\) when
\(L=1\); the boundary \(\sigma^4=h_L\) belongs to the second case.
% lean: CausalSmith.Stat.RdTruesideNoiseFrontier.variance_cost_envelope
\end{enumerate}
\end{proof}
